\chapter{Vitamin C Deficiency --- Scurvy}
\index{Vitamin C Deficiency --- Scurvy}
\label{sec:Vitamin C Deficiency}

\section{Overview}
Ascorbic acid (vitamin C) deficiency (scurvy) is a vitamin deficiency disorder that is rare in developed countries but occurs in individuals with severely restricted diets (alcoholics, elderly, psychiatric patients, food faddism), malabsorption, and in infants fed only cow's milk without supplementation. This metabolic disorder involves dysregulation of normal bodily processes, often requiring ongoing monitoring and management. It is increasingly common due to lifestyle and dietary factors. Metabolic disorders involve disruption of normal biochemical processes essential for health. These conditions may result from enzyme deficiencies, hormonal imbalances, or lifestyle factors. Many metabolic disorders are chronic and require ongoing management. Lifestyle modifications are often the cornerstone of treatment.
\section{Causes}
This condition happens because of deficiency of ascorbic acid, which is essential for collagen synthesis (hydroxylation of proline and lysine in procollagen), carnitine synthesis, neurotransmitter synthesis, and as an antioxidant. Metabolic disorders arise from dysregulation of normal biochemical pathways. This may result from enzyme deficiencies, hormonal imbalances, or lifestyle factors that overwhelm normal homeostatic mechanisms. The underlying mechanisms involve complex interactions between genetic predisposition and environmental factors. Disruption of normal physiological processes leads to the characteristic features of the condition. Multiple contributing factors may act synergistically to trigger or worsen the condition.
\section{Risk Factors}
\begin{itemize}[noitemsep]
  \item Genetic predisposition and family history
  \item Advancing age
  \item Lifestyle factors (diet, exercise, smoking, alcohol)
  \item Environmental exposures
  \item Underlying medical conditions
  \item Medication use
\end{itemize}
\section{Signs and Symptoms}

\subsection{Common symptoms}
\begin{itemize}[noitemsep]
  \item Clinical features reflect impaired collagen synthesis: perifollicular hemorrhages, tiny red or purple spots, ecchymoses, gingival abnormal increase in cells with bleeding, impaired wound healing, hyperkeratosis, coiled hairs (corkscrew hairs), and leg swelling
  \item Musculoskeletal manifestations include arthralgias, myalgias, and subperiosteal hemorrhages (particularly painful in infants presenting with pseudoparalysis)
  \item Systemic symptoms include fatigue, irritability, depression, and anemia.
  \item Pain or discomfort in the affected area
  \end{itemize}

\subsection{Emergency warning signs}
\begin{itemize}[noitemsep]
  \item Severe or worsening symptoms of vitamin c deficiency --- scurvy require immediate medical evaluation
\end{itemize}
\section{Progression and Stages}
The condition may progress through different stages of severity over time. Early stages often involve mild or subtle symptoms that may be overlooked. Moderate stages involve more noticeable symptoms affecting daily function. Advanced stages may involve significant impairment and complications.
\section{Complications}
\begin{itemize}[noitemsep]
  \item Disease progression leading to worsening symptoms
  \item Secondary infections or injuries
  \item Side effects of treatment
  \item Reduced quality of life
  \item Emotional and psychological impact
\end{itemize}
\section{Diagnosis}
\begin{itemize}[noitemsep]
  \item Doctors usually diagnose this based on your symptoms and a physical exam
  \item Serum ascorbic acid less than 0.2 mg/dL confirms deficiency.
  \item Comprehensive medical history and physical examination
  \item Laboratory tests to assess organ function
  \item Imaging studies to visualize affected structures
  \item Specialized testing depending on the affected system
  \item Consultation with relevant specialists
\end{itemize}
\section{Prevention}
\begin{itemize}[noitemsep]
  \item Maintain a healthy lifestyle with balanced diet and exercise
  \item Avoid known risk factors
  \item Attend regular medical check-ups for early detection
  \item Follow recommended screening guidelines
\end{itemize}
\section{Treatment Options}
\begin{itemize}[noitemsep]
  \item Treatment focuses on vitamin C 100--200 mg orally three times daily for 1 week, then 100 mg daily
  \item Infantile scurvy responds to 50 mg daily.
  \item Medications to manage symptoms and address underlying mechanisms
  \item Lifestyle modifications and self-care measures
  \item Physical therapy and rehabilitation
  \item Surgical intervention when indicated
  \item Regular monitoring and follow-up care
\end{itemize}
\section{Living with It}
\begin{itemize}[noitemsep]
  \item Follow your treatment plan as prescribed
  \item Maintain regular follow-up with your healthcare provider
  \item Make necessary lifestyle adjustments
  \item Seek support from family, friends, or support groups
  \item Stay informed about your condition
\end{itemize}
\section{Outlook}
\begin{itemize}[noitemsep]
  \item Most people recover well
  \item Clinical improvement occurs within days.
\end{itemize}
